% abstract, choose between abstract and summary

I recenti sviluppi tecnologici e produttivi registrati nel campo della robotica subacquea, con particolare riferimento ai veicoli subacquei autonomi (Autonomous Underwater Vehicles - AUV), hanno contribuito ad incrementare notevolmente le capacità operative di tali veicoli, consentendo loro di affrontare missioni di lunga durata e ampliandone le potenziali applicazioni nelle attività di sorveglianza, rilevamento e manutenzione. La diffusione su larga scala di questi sistemi richiede, tuttavia, un contenimento dei costi di realizzazione, perseguibile anche tramite l'impiego di sensori economicamente vantaggiosi. Questo compromesso, però, può comportare maggiori incertezze nelle misure acquisite, derivanti da fenomeni quali il rumore e gli errori sistematici. Esse, congiuntamente alle discrepanze tra il modello matematico e la dinamica reale del sistema, possono limitare significativamente la capacità delle tecniche model-based tradizionali, come quelle basate sui filtri di Kalman, di stimare accuratamente la posizione e lo stato del veicolo. Tali problematiche sono, inoltre, accentuate dall'impossibilità di fare affidamento su sistemi di posizionamento esterni, come il GNSS, in grado di correggere gli eventuali errori accumulati durante la navigazione subacquea. In questo lavoro di tesi, svolto nell'ambito di una collaborazione con l'azienda Modelway, si è cercato di superare queste limitazioni, esplorando le potenzialità dei più recenti algoritmi di apprendimento automatico, basati su un approccio puramente data-driven. Questi metodi, ascrivibili all'ambito dell'intelligenza artificiale, consentono, infatti, di apprendere relazioni complesse e non lineari direttamente dall'analisi dei dati acquisiti dai sensori di bordo, offrendo così una possibile alternativa, o integrazione, ai metodi di posizionamento classici. Dopo aver presentato i principali paradigmi dell'intelligenza artificiale, con particolare riferimento alle tecniche di machine learning e deep learning, si è, dunque, condotta un'analisi comparativa dei principali approcci proposti in letteratura, giungendo, infine, alla selezione dell'architettura NavNet come riferimento per questa tesi. Tale rete neurale, in particolare, si caratterizza per la presenza di più rami posti in parallelo, ognuno composto da un modulo LSTM e da un Simplified Attention Mechanism, i cui vettori di output sono sintetizzati, infine, da due Fully Connected Layer. Sulla base dei dati raccolti mediante l'utilizzo di un simulatore di un drone subacqueo commerciale, è stata, quindi, condotta un'analisi approfondita del modello selezionato, volta a valutarne le effettive prestazioni sullo specifico caso studio presentato e ad individuare possibili margini di miglioramento dell'architettura originaria.